% Source: OpenAI2026AreaLaw, 09-amplification.tex, lines 139--173,
% openai/math@adc7f1241b42e322a6451854ab7e4b4c146bf78a.
% ChannelWordObservables is checked at b9968a9236.
% These are finite chronological compositions and their append increments.
% No clock process or Poisson evolution estimate is asserted here.

\section{Chronological finite channel words}
\label{sec:al_channel_words}

Let $I,J$ be finite sets, let $q\in\N$, and put
$\mathcal H=(\C^q)^{\otimes I}$ and $\mathcal K=\C^J$.
The spectator dimension may be zero. Let $F$ be a set of event labels,
and write $\mathcal E_i$ for the full spectator root map associated
with the physical matrix $k_i$, as in
Definition~\ref{def:al_local_root_channel}.
All input operators are arbitrary complex matrices on
$\mathcal H\otimes\mathcal K$, with the Euclidean operator norm.
Write $wi$ for the word obtained by appending a single label $i$ to
the right of a finite word $w$.

\begin{definition}[Observable after a chronological word]
    \label{def:al_channel_word}
    \uses{def:al_local_root_channel}
    \lean{TNLean.PEPS.AreaLaw.spectatorRootChannelWord}
    \leanok
    For a nonempty word $w=(i_0,\ldots,i_{m-1})$, define
    \begin{align}
        \label{eq:al_word_composition}
        \mathcal E_w(B)
         & =(\mathcal E_{i_{m-1}}\circ\cdots\circ\mathcal E_{i_0})(B).
    \end{align}
    For the empty word $\varnothing$, take
    $\mathcal E_{\varnothing}=\mathrm{id}$.
    Thus the first event acts first. Given an initial operator $B_0$,
    write $B_w=\mathcal E_w(B_0)$.
\end{definition}

\begin{theorem}[Empty words and appending an event]
    \label{thm:al_channel_word_append}
    \uses{def:al_channel_word}
    \lean{TNLean.PEPS.AreaLaw.spectatorRootChannelWord_nil,
        TNLean.PEPS.AreaLaw.spectatorRootChannelWord_append_singleton}
    \leanok
    For every operator $B$, every finite word $w$, and every label $i$,
    \begin{align}
        \label{eq:al_word_empty}
        \mathcal E_{\varnothing}(B) & =B,
    \end{align}
    \begin{align}
        \label{eq:al_word_append}
        \mathcal E_{wi}(B) & =\mathcal E_i(\mathcal E_w(B)).
    \end{align}
    In particular, $B_{\varnothing}=B_0$ and $B_{wi}=\mathcal E_i(B_w)$.
\end{theorem}
\begin{proof}
    \leanok
    \uses{def:al_channel_word}
    The empty composition is the identity. In the composition for
    $wi$, the last map is $\mathcal E_i$, after every map indexed by
    $w$, which gives~(\ref{eq:al_word_append}).
\end{proof}

\begin{theorem}[Norm contraction along a finite word]
    \label{thm:al_channel_word_norm}
    \uses{def:al_channel_word}
    \lean{TNLean.PEPS.AreaLaw.norm_spectatorRootChannelWord_le}
    \leanok
    Suppose $0\le k_i\le\mathbb 1$ for every $i\in F$.
    Then $\|\mathcal E_w(B)\|\le\|B\|$ for every finite word $w$
    and every operator $B$.
\end{theorem}
\begin{proof}
    \leanok
    \uses{def:al_channel_word,thm:al_root_channel_contraction}
    Induct on the word, allowing an arbitrary input operator.
    The empty word gives equality. If the first label is $i$ and
    the remaining word is $u$, then
    \begin{align}
        \label{eq:al_word_norm_step}
        \|\mathcal E_{iu}(B)\|
         & =\|\mathcal E_u(\mathcal E_i(B))\|
        \le\|\mathcal E_i(B)\|\le\|B\|.
    \end{align}
    Here the first inequality is the induction hypothesis at input
    $\mathcal E_i(B)$ and the second is
    Theorem~\ref{thm:al_root_channel_contraction}.
\end{proof}

\begin{theorem}[Oscillation bounded by the initial norm]
    \label{thm:al_channel_word_oscillation}
    \uses{def:al_channel_word,def:qc_site_oscillation}
    \lean{TNLean.PEPS.AreaLaw.siteOscillation_spectatorRootChannelWord_bounds}
    \leanok
    Suppose $0\le k_i\le\mathbb 1$ for every $i\in F$.
    For every finite word $w$, site $y\in I$, and initial operator
    $B_0$, one has $0\le d_y(B_w)\le2\|B_0\|$.
\end{theorem}
\begin{proof}
    \leanok
    \uses{thm:al_channel_word_norm,thm:qc_site_oscillation_bounds}
    The oscillation is nonnegative, and
    $d_y(B_w)\le2\|B_w\|\le2\|B_0\|$ by
    Theorems~\ref{thm:qc_site_oscillation_bounds} and
    \ref{thm:al_channel_word_norm}.
\end{proof}

At the empty word, the initial support bound is
Theorem~\ref{thm:al_spectator_support_indicator}.
If $q>0$ and every physical matrix block of $B_0$ is supported on
$S\subseteq I$, then~(\ref{eq:al_word_empty}) gives
\begin{align}
    \label{eq:al_word_initial_support}
    d_y(B_{\varnothing}) & =d_y(B_0)\le2\|B_0\|\mathbf1_S(y).
\end{align}

\begin{theorem}[Summed append increments of the current observable]
    \label{thm:al_channel_word_kernel_increment}
    \uses{def:al_channel_word,def:qc_site_oscillation,
        def:al_graph_event_kernel,def:qc_local_circuit}
    \lean{TNLean.PEPS.AreaLaw.sum_siteOscillation_spectatorRootChannelWord_append_sub_le_graphChannelEventKernel}
    \leanok
    Suppose $q>0$ and $F$ is finite. Let $G$ be a simple graph on
    $I$, and choose anchors $a_i\in I$ for $i\in F$.
    Assume $0\le k_i\le\mathbb 1$ and that $k_i$ is supported on
    the connected component of $a_i$ for every $i$.
    Let $C\ge0$, $c>0$, and $0<\alpha\le1$.
    Fix a cutoff $N\in\N$ with $N\ge R_{a_i}$ for every $i$,
    where $R_{a_i}$ is the maximum distance from $a_i$ within its
    component. Suppose that, for every $i\in F$ and every integer
    $l$ with $0\le l\le N$,
    \begin{align}
        \label{eq:al_word_localization_error}
        \|k_i-E_{N_l(a_i)}(k_i)\| & \le C\exp(-cl^\alpha).
    \end{align}
    Use the kernel~(\ref{eq:al_propagation_event_kernel}) with
    $D=2(2+8\sqrt C\exp(c/2))$ and $b=c/2$.
    Then, for every finite word $w$, every $y\in I$, and every
    initial operator $B_0$,
    \begin{align}
        \label{eq:al_word_kernel_increment}
        \sum_{i\in F}\bigl(d_y(B_{wi})-d_y(B_w)\bigr)
         & \le\sum_{z\in I}\mathsf A^{(N)}_{yz}d_z(B_w).
    \end{align}
    Every oscillation on the right is evaluated at the same current
    observable $B_w$. Labels with the same anchor remain distinct
    summands. The estimate holds for arbitrary finite spectators
    and arbitrary complex initial operators $B_0$.
\end{theorem}
\begin{proof}
    \leanok
    \uses{thm:al_actual_channel_kernel_domination,thm:al_channel_word_append}
    Apply Theorem~\ref{thm:al_actual_channel_kernel_domination} to
    the operator $B_w$. For each event $i$, replace
    $\mathcal E_i(B_w)$ by $B_{wi}$ using~(\ref{eq:al_word_append}).
    The resulting inequality is~(\ref{eq:al_word_kernel_increment}).
\end{proof}
